% =====================================================================
%  Legend of the Technomancers — OPEN QUESTIONS
%  Every design note, TODO, and unresolved fork, pulled OUT of the
%  rulebook so the book reads as a finished product.
%  Build:  tectonic -X compile open-questions.tex
% =====================================================================
\documentclass[11pt,oneside]{memoir}
\usepackage{atrpg}
\geometry{a4paper, inner=20mm, outer=20mm, top=20mm, bottom=22mm}

% \leavevmode first: ragged2e is loaded with [document] and hooks
% \everyselectfont, so a bare font switch in vertical mode drops a line.
\newcommand{\thc}[1]{\leavevmode{\hdfont\bfseries\color{white}#1}}
\newcommand{\where}[1]{\par\smallskip{\small\itshape\color{atrAccent2}Was in: #1}\par}

\begin{document}
\color{atrInk}

% ------------------------------------------------------------ Title page
\begin{titlingpage}
\begin{tikzpicture}[remember picture, overlay]
  \fill[atrInk] (current page.north west) rectangle ([yshift=-104mm]current page.north east);
  \node[anchor=north, align=center, inner sep=0] at ([yshift=-30mm]current page.north){%
    \thc{\fontsize{34}{38}\selectfont LEGEND OF THE}\\[2mm]
    \thc{\fontsize{34}{38}\selectfont TECHNOMANCERS}\\[8mm]
    {\color{atrAccent2}\rule{0.50\paperwidth}{1.2pt}}\\[8mm]
    {\hdfont\bfseries\color{atrAccent2}\fontsize{28}{32}\selectfont OPEN QUESTIONS}\\[7mm]
    {\hdfont\large\itshape\color{white}Everything Still To Decide}%
  };
  \node[anchor=south, align=center, inner sep=0] at ([yshift=26mm]current page.south){%
    {\itshape None of this belongs in the rulebook.}\\
    {\itshape It lives here until it is decided.}%
  };
\end{tikzpicture}
\end{titlingpage}

\frontmatter
\tableofcontents

\mainmatter

% =====================================================================
\chapter{How to Use This File}
% =====================================================================

\lettrine{T}{his file} is the holding pen. The rulebook is meant to read
like a finished game --- a reader wants the rules, not a tour of how we got
to them --- so every design note, open question, and \emph{to do} has been
pulled out of it and collected here.

\begin{statblock}[The Rule]
\begin{tabularx}{\linewidth}{@{}p{4.2cm} Y@{}}
\keyword{In the book} & Rules as they currently stand, stated plainly and without apology. Reader-facing guidance --- including advice to the GM --- is fine.\\[3pt]
\keyword{In this file} & Anything unresolved, any rationale for a choice, any ``working idea'', and any note that the text is incomplete.\\[3pt]
\keyword{When decided} & Write the rule into the book in its own voice, and delete the entry here.\\
\end{tabularx}
\end{statblock}

Each entry records where it used to sit, so it is easy to put the answer
back in the right place.

% =====================================================================
\chapter{Settled, For the Record}
% =====================================================================

These were decided on 2026-09-21 and are already written into the book.
They are here only because the \emph{reasoning} is worth keeping, and the
reasoning is not the reader's business.

\section{Why There Is No Magic Stat}
There is deliberately no Magic or Arcane stat. If there were, a character
could pour everything into one hugely versatile number and dominate every
kind of action at once. Hanging magic off the same three stats keeps tech,
diplomacy, mundane expertise, and sorcery equally powerful --- and it makes
\emph{how} you work magic say something about who you are.
\where{How to Play}

\section{Why Three Stats, Not Four}
Body, Mind, and Face replaced Strength, Agility, Charisma, and Intellect
because the same three distinctions apply whether an action is mundane or
magical, which is what lets one engine cover both.

% =====================================================================
\chapter{The Core Engine}
% =====================================================================

\section{``Body'' Means Two Things}
\keyword{Naming clash to resolve.} \emph{Body} is currently both one of the
three Stats \emph{and} one of the four creation columns, so ``your Body rank
sets your Body stat'' reads badly.

\begin{statblock}[Options]
\begin{tabularx}{\linewidth}{@{}p{4.2cm} Y@{}}
\keyword{Rename the stat} & \emph{Physique}, \emph{Brawn}, \emph{Form}.\\[3pt]
\keyword{Rename the column} & \emph{Flesh}, \emph{Self}, \emph{Shape}.\\[3pt]
\keyword{Accept it} & Lean on context; the two are rarely adjacent.\\
\end{tabularx}
\end{statblock}
\where{How to Play}

\section{Higher of the Two}
\textbf{To do:} confirm that \emph{take the higher} of character stat and
item tier is balanced, rather than always using the object's tier, or
averaging the two. This is the working rule and it wants play-testing.
\where{How to Play --- Stats from Objects}

\section{Guidelines for Setting Difficulties}
\textbf{To do:} the GM sets a target number and there is currently no
guidance on what target to pick. Note that with success on 4--6, a target of
half your pool is always a coin flip, and the useful band is narrow --- 3 to
5 successes. Every number the GM has to invent is friction.

\section{Hit Points}
\textbf{To do:} decide whether durability is the \keyword{Body} stat
directly, a pool derived from it, or something tied to all three stats.
\where{Body}

% =====================================================================
\chapter{Body}
% =====================================================================

\section{Body Points}
\textbf{To do:} decide how many buying points a Body rank grants, and the
stat cap at each rank. Stats run \dice{1}--\dice{5}; ranks \rating{3} and
\rating{4} should also grant qualitatively different traits --- size, a
natural weapon, supernatural resilience.
\where{Body --- At Character Creation}

\section{Body Tags Versus Modifications}
\keyword{Open direction:} rather than a fixed Modification list, Body could
grant \keyword{tags} the way Magic does --- a shared set of common tags
(\emph{Attractive}, \emph{Fast}, \emph{Armoured}) plus world-specific ones,
so a demon body might carry \emph{Fire}. Tags would then feed rolls from
four sources alike: body, stuff, skills, and magic.

There is a concrete reason to prefer it: collapsing four stats into three
left \emph{Extra Strong} and \emph{Fleet of Foot} both granting \dice{+2}
Body, which are now the same option wearing two names. Tags would carry that
distinction where a stat bonus no longer can.
\where{Body --- At Character Creation}

\section{Mobility}
\textbf{To do:} write up special movement --- flight, shadow-walking,
web-walking, supernatural speed. These can come from a body, from magic, or
from an item, and they are a large part of what makes a rank \rating{3} or
\rating{4} body feel strange.
\where{Body --- Your Modifications}

% =====================================================================
\chapter{Stuff}
% =====================================================================

\section{Purchasing}
\textbf{To do:} decide how many purchasing points a character gets, what a
point buys at each rank, and the upper limits. The ambition is a table of
costs like the one in the \emph{D\&D} Player's Handbook.
\where{Stuff --- At Character Creation}

\section{Vehicles in Play}
\textbf{To do:} decide what a vehicle does mechanically --- chase rules,
how it soaks damage, how many it carries, and whether higher tiers grant
dice on the relevant rolls.
\where{Stuff --- Vehicles}

\section{Damage and Armour}
\textbf{To do:} settle how damage works. Working idea: each rank of armour
reduces incoming damage by one point.
\where{Stuff --- Everything Else}

\section{Legendary Weapons}
\textbf{To do:} a tier-4 weapon is currently described as ``a weapon out of
story'' because nobody has decided what it actually does. It needs a real
mechanical hook.

% =====================================================================
\chapter{Skills}
% =====================================================================

\section{Skill Points}
\textbf{To do:} pin down how many skill points each priority grants and
whether any single skill is capped. Working idea: \rating{1} = 4 points,
\rating{2} = 8, \rating{3} = 14, \rating{4} = 20.
\where{Skills --- At Character Creation}

\section{Does the Stat Column Earn Its Keep?}
With only three Stats, each skill's listed pairing is nearly implied by its
category: Conflict and Movement are \keyword{Body}, Social is
\keyword{Face}, Knowledge and Technical are \keyword{Mind}. If the Stat is
chosen by the fiction anyway, the column may be redundant --- worth cutting,
or demoting to a parenthetical hint.
\where{Skills --- Your Skills}

\section{Genre-Specific Skills}
\textbf{To do:} the list in the book is the \emph{common} set, meant to work
in any realm. Genre-specific skills (\emph{Cyberdeck Operation},
\emph{Starship Gunnery}, \emph{Heraldry}) want a separate treatment ---
working idea: you hold a fixed number of them, and decide what they are as
they come up in play.
\where{Skills --- Your Skills}

% =====================================================================
\chapter{Magic and Tags}
% =====================================================================

\section{Tag Breadth}
Broad tags are mechanically stronger, because they apply more often. The
defences are a curated list, worked examples of acceptable scope, and the GM
veto: \emph{Control} and \emph{Animals} are fine, \emph{Rats} and
\emph{Underground} are good specific tags, and \emph{Adventuring} is too
broad and should be refused.

\textbf{Open:} should a narrow tag grant an explicit bonus, or is the fact
that it applies more rarely already payment enough?

\begin{notebox}[Worth Promoting]
The curated-list-plus-GM-veto guidance is real, usable rules text. When the
narrow-tag question is settled, this should go into the book as guidance ---
stated as rules, not as a note about rules.
\end{notebox}
\where{Magic --- How Magic Works}

\section{How Many Tags May Contribute}
\textbf{Open:} the Magic rank caps applicable tags at 2/3/5. Once tags also
come from body, stuff, and skills, is there a single cap across all sources,
or one per source?

\section{Do Tags Gate, or Gate and Count?}
\textbf{Open fork.} Andrew wants the columns to work as a permissions system
--- you can only do a thing if you have the capability for it. If tags both
\emph{gate} and \emph{add dice}, then a character with zero applicable tags
cannot act at all. That is a hard wall, and very much a permissions game,
but it needs a deliberate yes or no.

\section{Genre Swap Semantics}
\textbf{Open fork.} When a character crosses realms, do their tags
\emph{persist and get re-described} (a dragon's \emph{fire, flight,
armoured} become a war-mech's, same tags), or do they get \emph{re-picked}
(\emph{fire/flight/claws} traded for \emph{steam/servos/plating})?

Re-picking is more evocative --- it is what makes ``a fiery dragon becomes a
steaming robot'' land --- but it is a character rebuild at every border, and
it invites optimising for the adventure you can already see. Middle path:
carry tags over re-skinned by default, and allow swapping up to \emph{rank}
of them on arrival.

% =====================================================================
\chapter{Advancement and the Sheet}
% =====================================================================

\section{Points Conversion}
\textbf{To do:} figure out the points-conversion rules --- how Change Points
and Persistence Points are earned, and their relative costs.
\where{Character Advancement}

\section{Sheet Layout}
\textbf{To do:} design the character sheet itself --- a one-page, printable
layout, themed to match the book. It needs \emph{check-off} sections for the
fixed options (the base Magic tags, the common skills) and \emph{free-entry}
sections for the world-specific things a traveller picks up in each realm.
\where{The Character Sheet}

\section{Visual Language}
\textbf{To do:} two shapes carry the whole sheet.
\begin{itemize}
  \item A \keyword{Stat} is an icon plus a number in a \emph{rectangle} --- a muscle for Body, and so on.
  \item A \keyword{tag} is a word in an \emph{oval}, carrying a small icon for its provenance: magic, body, skill, or item.
\end{itemize}
Once tags come from all four sources, that provenance icon is what keeps the
sheet readable at a glance.
\where{The Character Sheet}

\section{Example Characters}
\textbf{To do:} write a set of finished example characters that lean hard
into recognisable tropes, so a new player can see where the numbers land
before building their own.
\where{The Character Sheet}

% =====================================================================
\chapter{Unwritten Content}
% =====================================================================

These are not design questions --- the design is known, the prose simply is
not written. The book currently carries \verb|\lipsum| placeholders in these
places.

\begin{statblock}[Still To Write]
\begin{tabularx}{\linewidth}{@{}p{4.6cm} Y@{}}
\keyword{Part II: Playing} & Running conflicts, session pacing, travel between realms, the GM's role, and mid-campaign rewards. Currently a stub.\\[3pt]
\keyword{High Fantasy reskins} & The realm's take on Weapons, Armour, Lifestyle, and Locations, then on Body, Skills, and Magic.\\[3pt]
\keyword{Opening fiction} & Every chapter's read-aloud box is placeholder text.\\[3pt]
\keyword{The pitch} & The intro never says what the game is or what its tone is.\\[3pt]
\keyword{Per-realm bestiaries} & Floated as the way to handle Body/Stuff interacting with Magic/Skills, which is genre-specific and hard to spec generically. Also supplies each realm's tag vocabulary.\\[3pt]
\keyword{World-gen sliders} & Mentioned in the notes, not designed.\\
\end{tabularx}
\end{statblock}

\end{document}
